% GENERATED FILE. Edit canonical lessons, then run npm run generate:books.
% canonical-source-hash: fnv1a64:2ded2530e3de15a4
% canonical-lessons: AR-C371-wasila, AR-C371-saa, AR-C371-anjaza, AR-C371-sahama, AR-C371-hafiz

\chapter{Aims and means}
\label{ch:ar-aims-and-means}
\hlchaptermodality{371}

\begin{chapteropening}
\textbf{By the end of this chapter:} \emph{I can talk about a means and an incentive, and say that someone strove, accomplished or contributed.}

Complete the last lesson of chapter 371: I can talk about a means and an incentive, and say that someone strove, accomplished or contributed.
\end{chapteropening}

\section[\texorpdfstring{\ar{وسيلة} (wasīla)}{وسيلة (wasīla)}]{\ar{وسيلة} (wasīla) — a means}

\label{lesson:AR-C371-wasila}



\begin{quote}
Before the new one: say the Arabic for to exaggerate, then the Arabic for to generalize.
\end{quote}

\subsection*{You'll want to know: \ar{وسيلة}}
\textbf{\ar{وسيلة}} — \emph{wasīla} — \textquotedblleft{}a means\textquotedblright{}.

\begin{quote}
\ar{الحافلة} \ar{وسيلة} \ar{رخيصة} \ar{للسفر}.
\end{quote}

\begin{quote}
\emph{al-ḥāfila wasīla rakhīṣa li-s-safar.}
\end{quote}

\begin{quote}
The bus is a cheap means of travel.
\end{quote}

\begin{cousinweb}
\textbf{\ar{وسيلة}} (\emph{wasīla}) is built on the native Arabic root \emph{w-s-l}, \textquotedblleft{}to seek to draw near\textquotedblright{} (Lane). A \emph{wasīla} is what brings you near to a goal: a means.
\end{cousinweb}

\begin{grammarlens}[title={What you've built}]
One more word for giving reasons and opinions.
\end{grammarlens}

\subsection*{Your turn}
\begin{itemize}
\raggedright
  \item \emph{Say it:} \emph{wasīla}
  \item \emph{Say it:} \emph{wasīla}, once more
  \item \emph{Say it:} \emph{al-ḥāfila wasīla rakhīṣa li-s-safar.} — The bus is a cheap means of travel.
  \item \emph{Recall:} say the Arabic for to exaggerate, then the Arabic for to generalize, then say \emph{wasīla} again
\end{itemize}

\begin{tcolorbox}[breakable,colback=teal!4,colframe=teal!35!black,title={Before you move on}]
What does \ar{وسيلة} mean? (\textquotedblleft{}A means\textquotedblright{}.) Say it once more.
\end{tcolorbox}
\section[\texorpdfstring{\ar{سعى} (saʿā)}{سعى (saʿā)}]{\ar{سعى} (saʿā) — to strive}

\label{lesson:AR-C371-saa}



\begin{quote}
Before the new one: say the Arabic for to generalize, then the Arabic for a means.
\end{quote}

\subsection*{You'll want to know: \ar{سعى}}
\textbf{\ar{سعى}} — \emph{saʿā} — \textquotedblleft{}to strive\textquotedblright{}.

\begin{quote}
\ar{سعى} \ar{الشاب} \ar{إلى} \ar{عمل} \ar{أفضل}.
\end{quote}

\begin{quote}
\emph{saʿā ash-shābb ilā ʿamal afḍal.}
\end{quote}

\begin{quote}
The young man strove for a better job.
\end{quote}

Present: \ar{يسعى}, \emph{yasʿā}.

It takes \ar{إلى}: \ar{سعى} \ar{إلى} ….

\begin{cousinweb}
\textbf{\ar{سعى}} (\emph{saʿā}) is built on the native Arabic root \emph{s-ʿ-y}, usually traced to the sense \textquotedblleft{}to run, to strive\textquotedblright{}. \emph{Saʿā} is a Form I verb with a weak last letter, so its present is \emph{yasʿā}.
\end{cousinweb}

\begin{grammarlens}[title={What you've built}]
One more word for giving reasons and opinions.
\end{grammarlens}

\subsection*{Your turn}
\begin{itemize}
\raggedright
  \item \emph{Say it:} \emph{saʿā}
  \item \emph{Say it:} \emph{saʿā}, once more
  \item \emph{Say it:} \emph{saʿā ash-shābb ilā ʿamal afḍal.} — The young man strove for a better job.
  \item \emph{Recall:} say the Arabic for to generalize, then the Arabic for a means, then say \emph{saʿā} again
\end{itemize}

\begin{tcolorbox}[breakable,colback=teal!4,colframe=teal!35!black,title={Before you move on}]
What does \ar{سعى} mean? (\textquotedblleft{}To strive\textquotedblright{}.) Say it once more.
\end{tcolorbox}
\section[\texorpdfstring{\ar{أنجز} (anjaza)}{أنجز (anjaza)}]{\ar{أنجز} (anjaza) — to accomplish}

\label{lesson:AR-C371-anjaza}



\begin{quote}
Before the new one: say the Arabic for a means, then the Arabic for to strive.
\end{quote}

\subsection*{You'll want to know: \ar{أنجز}}
\textbf{\ar{أنجز}} — \emph{anjaza} — \textquotedblleft{}to accomplish\textquotedblright{}.

\begin{quote}
\ar{أنجز} \ar{العمال} \ar{البناء} \ar{في} \ar{سنة}.
\end{quote}

\begin{quote}
\emph{anjaza al-ʿummāl al-bināʾ fī sana.}
\end{quote}

\begin{quote}
The workers completed the building in a year.
\end{quote}

Present: \ar{ينجز}, \emph{yunjiz}.

\begin{cousinweb}
\textbf{\ar{أنجز}} (\emph{anjaza}) is built on the native Arabic root \emph{n-j-z}, \textquotedblleft{}to be accomplished, finished\textquotedblright{} (Lane). It is Form IV, to bring to completion: to accomplish.
\end{cousinweb}

\begin{grammarlens}[title={What you've built}]
One more word for giving reasons and opinions.
\end{grammarlens}

\subsection*{Your turn}
\begin{itemize}
\raggedright
  \item \emph{Say it:} \emph{anjaza}
  \item \emph{Say it:} \emph{anjaza}, once more
  \item \emph{Say it:} \emph{anjaza al-ʿummāl al-bināʾ fī sana.} — The workers completed the building in a year.
  \item \emph{Recall:} say the Arabic for a means, then the Arabic for to strive, then say \emph{anjaza} again
\end{itemize}

\begin{tcolorbox}[breakable,colback=teal!4,colframe=teal!35!black,title={Before you move on}]
What does \ar{أنجز} mean? (\textquotedblleft{}To accomplish\textquotedblright{}.) Say it once more.
\end{tcolorbox}
\section[\texorpdfstring{\ar{ساهم} (sāhama)}{ساهم (sāhama)}]{\ar{ساهم} (sāhama) — to contribute}

\label{lesson:AR-C371-sahama}



\begin{quote}
Before the new one: say the Arabic for to strive, then the Arabic for to accomplish.
\end{quote}

\subsection*{You'll want to know: \ar{ساهم}}
\textbf{\ar{ساهم}} — \emph{sāhama} — \textquotedblleft{}to contribute\textquotedblright{}.

\begin{quote}
\ar{ساهم} \ar{الجيران} \ar{في} \ar{تنظيف} \ar{الشارع}.
\end{quote}

\begin{quote}
\emph{sāhama al-jīrān fī tanẓīf ash-shāriʿ.}
\end{quote}

\begin{quote}
The neighbours contributed to cleaning the street.
\end{quote}

Present: \ar{يساهم}, \emph{yusāhim}.

It takes \ar{في}: \ar{ساهم} \ar{في} ….

\begin{cousinweb}
\textbf{\ar{ساهم}} (\emph{sāhama}) is built on the native Arabic root \emph{s-h-m}, \textquotedblleft{}an arrow; a share\textquotedblright{} (Lane). It is Form III, to take a share with others: to contribute.
\end{cousinweb}

\begin{grammarlens}[title={What you've built}]
One more word for giving reasons and opinions.
\end{grammarlens}

\subsection*{Your turn}
\begin{itemize}
\raggedright
  \item \emph{Say it:} \emph{sāhama}
  \item \emph{Say it:} \emph{sāhama}, once more
  \item \emph{Say it:} \emph{sāhama al-jīrān fī tanẓīf ash-shāriʿ.} — The neighbours contributed to cleaning the street.
  \item \emph{Recall:} say the Arabic for to strive, then the Arabic for to accomplish, then say \emph{sāhama} again
\end{itemize}

\begin{tcolorbox}[breakable,colback=teal!4,colframe=teal!35!black,title={Before you move on}]
What does \ar{ساهم} mean? (\textquotedblleft{}To contribute\textquotedblright{}.) Say it once more.
\end{tcolorbox}
\section[\texorpdfstring{\ar{حافز} (ḥāfiz)}{حافز (ḥāfiz)}]{\ar{حافز} (ḥāfiz) — an incentive}

\label{lesson:AR-C371-hafiz}



\begin{quote}
Before the new one: say the Arabic for to accomplish, then the Arabic for to contribute.
\end{quote}

\subsection*{You'll want to know: \ar{حافز}}
\textbf{\ar{حافز}} — \emph{ḥāfiz} — \textquotedblleft{}an incentive\textquotedblright{}.

\begin{quote}
\ar{كان} \ar{الراتب} \ar{الجيد} \ar{حافزا} \ar{للعمل}.
\end{quote}

\begin{quote}
\emph{kāna ar-rātib al-jayyid ḥāfizan li-l-ʿamal.}
\end{quote}

\begin{quote}
The good salary was an incentive to work.
\end{quote}

Plural: \ar{حوافز}, \emph{ḥawāfiz}.

\begin{cousinweb}
\textbf{\ar{حافز}} (\emph{ḥāfiz}) is built on the native Arabic root \emph{ḥ-f-z}, \textquotedblleft{}to urge on from behind\textquotedblright{} (Lane). \emph{Ḥāfiz} is the active participle: what urges you on, an incentive.
\end{cousinweb}

\begin{grammarlens}[title={What you've built}]
One more word for giving reasons and opinions.
\end{grammarlens}

\subsection*{Your turn}
\begin{itemize}
\raggedright
  \item \emph{Say it:} \emph{ḥāfiz}
  \item \emph{Say it:} \emph{ḥāfiz}, once more
  \item \emph{Say it:} \emph{kāna ar-rātib al-jayyid ḥāfizan li-l-ʿamal.} — The good salary was an incentive to work.
  \item \emph{Recall:} say the Arabic for to accomplish, then the Arabic for to contribute, then say \emph{ḥāfiz} again
\end{itemize}

\begin{tcolorbox}[breakable,colback=teal!4,colframe=teal!35!black,title={Before you move on}]
What does \ar{حافز} mean? (\textquotedblleft{}An incentive\textquotedblright{}.) Say it once more.
\end{tcolorbox}
